\documentclass[11pt]{book}
\usepackage[paperwidth=145mm,paperheight=200mm,top=18mm,bottom=20mm,inner=20mm,outer=16mm]{geometry}
\usepackage{brumfiel}

% bookindex.tex -- GENERATED by tools/build_index.py; do not hand-edit.
% The entries and their nesting are the 1960 book's own index; the page
% numbers are this edition's, resolved by locating each entry's words inside
% the chapter the original reference pointed into.

% ---- local macros (hoist into book preamble) ----
\newcommand{\idxmain}[2]{\par\noindent\hangindent=1.2em #1\quad #2\par}
\newcommand{\idxsub}[2]{\par\noindent\hspace*{1.2em}\hangindent=2.4em #1\quad #2\par}
% ---- end local macros ----

\begin{document}
\setcounter{chapter}{19}

\chapter*{\raggedleft\Huge INDEX}
\addcontentsline{toc}{chapter}{Index}
\markboth{INDEX}{}
\vspace{-1em}\noindent\rule{\linewidth}{0.6pt}\vspace{1em}

\begin{multicols}{2}
\small\raggedright
\idxmain{Abscissa}{388}
\idxmain{Adams, John, C.}{413}
\idxmain{Alexander the Great}{6}
\idxmain{Algebra review}{59, 89, 104, 119, 168, 195, 222}
\idxmain{Analytic geometry}{384}
\idxmain{Angle(s)}{}
\idxsub{acute}{160}
\idxsub{adjacent}{79}
\idxsub{alternate exterior}{208}
\idxsub{alternate interior}{208}
\idxsub{congruence of}{121}
\idxsub{complementary}{210}
\idxsub{corresponding}{208}
\idxsub{definition of}{77}
\idxsub{definition of degree measure of}{331}
\idxsub{definition of "greater than" and "less than" for}{156}
\idxsub{definition of one degree}{331}
\idxsub{definition of right}{127}
\idxsub{existence of right}{152}
\idxsub{exterior of}{81}
\idxsub{included}{129}
\idxsub{interior of}{81}
\idxsub{measurement of}{330}
\idxsub{notation for indicating}{77}
\idxsub{obtuse}{160}
\idxsub{sides of}{77}
\idxsub{straight}{77}
\idxsub{sum of two}{212}
\idxsub{supplementary}{126}
\idxsub{supplementary and adjacent}{78}
\idxsub{vertex of}{77}
\idxsub{vertical}{78}
\idxmain{Apollonius}{6}
\idxmain{Arc}{317}
\idxmain{Archimedes}{8, 111}
\idxsub{Archimedes' axiom}{112}
\idxsub{Archimedes' integer}{111}
\idxmain{Area}{}
\idxsub{of a circle}{313}
\idxsub{determined by counting}{273}
\idxsub{as a limit}{275}
\idxsub{of a polyhedron}{375}
\idxsub{of a rectangle}{276}
\idxsub{of a sphere}{380}
\idxsub{of a triangle}{280}
\idxmain{Aristotle}{5, 411}
\idxmain{Assumptions}{63}
\idxmain{Axioms (postulates, assumptions)}{62}
\idxmain{Axis(es)}{}
\idxsub{of abscissas}{388}
\idxsub{coordinate}{388}
\idxsub{of ordinates}{388}
\idxmain{Bell, E. T.}{386, 436}
\idxmain{Bisector}{}
\idxsub{of an angle}{128}
\idxsub{perpendicular}{160}
\idxmain{Bolyai, J.}{7, 202}
\idxmain{Circle(s)}{}
\idxsub{center of}{180}
\idxsub{central angle of}{181}
\idxsub{chord of}{181}
\idxsub{circumscribing a polygon}{292}
\idxsub{definition of}{180}
\idxsub{definition of area of}{313}
\idxsub{definition of circumference of}{306}
\idxsub{diameter of}{180}
\idxsub{exterior of}{180}
\idxsub{inscribed angle of}{181}
\idxsub{inscribed in a polygon}{296}
\idxsub{interior of}{180}
\idxsub{radius of}{180}
\idxsub{secant of}{181}
\idxsub{tangent to}{294}
\idxmain{Circular arc(s)}{}
\idxsub{adjacent}{323}
\idxsub{definition of}{318}
\idxsub{definition of degree measure of}{331}
\idxsub{definition of semicircular arcs}{318}
\idxsub{midpoint of}{320}
\idxsub{subtended by a central angle}{318}
\idxmain{Commensurable segments}{242}
\idxmain{Completeness postulate}{114}
\idxmain{Congruence of angles}{}
\idxsub{postulate for}{121}
\idxsub{reflexive, symmetric, and transitive properties}{123}
\idxmain{Congruence of segments}{}
\idxsub{postulates for}{93}
\idxsub{reflexive, symmetric, and transitive properties}{94}
\idxsub{undefined relation}{98}
\idxmain{Congruence of triangles}{}
\idxsub{angle-side-angle theorem}{143}
\idxsub{definition of}{135}
\idxsub{side-angle-side postulate}{128}
\idxsub{side-angle-side theorem}{136}
\idxsub{side-side-side theorem}{144}
\idxmain{Converse of an implication}{44}
\idxmain{Coordinate(s)}{}
\idxsub{abscissa of a point}{392}
\idxsub{ordinate of a point}{392}
\idxsub{origin of}{388}
\idxsub{of a point in the plane}{392}
\idxsub{of a point on a line}{386}
\idxmain{Decagon}{216}
\idxmain{Definitions}{}
\idxsub{importance of}{103}
\idxsub{in algebra}{89, 104}
\idxmain{Degree measure}{}
\idxsub{of angles}{331}
\idxsub{of arcs}{334}
\idxmain{Descartes, René}{384}
\idxsub{his Discourse on Method}{384}
\idxsub{his Geometry}{384}
\idxmain{Diameter}{180}
\idxmain{Direction, rays with the same}{199}
\idxmain{Distance between points}{401}
\idxmain{Division}{}
\idxsub{definition in algebra}{89}
\idxsub{of segment into congruent parts}{235}
\idxmain{Equation(s)}{}
\idxsub{of a circle}{406}
\idxsub{equivalent}{399}
\idxsub{linear}{399}
\idxmain{Euclid}{7, 14, 173}
\idxsub{his Elements}{6, 14, 172}
\idxsub{his parallel postulate}{}
\idxmain{Eudoxus}{5, 112, 243}
\idxmain{Figures, relation between figures and geometry}{65}
\idxmain{Galileo}{411}
\idxmain{Gauss, K. F.}{7, 203}
\idxmain{Geometric figure}{11}
\idxmain{Geometry}{}
\idxsub{analytic}{385}
\idxsub{cartesian}{385}
\idxsub{coordinate}{385}
\idxsub{origin of word}{3}
\idxsub{reasons for study}{1}
\idxsub{of space}{360}
\idxmain{Greek letters}{31}
\idxmain{Heptagon}{216}
\idxmain{Heron's formula}{283}
\idxmain{Hexagon}{218}
\idxmain{Hilbert, David}{7, 95, 172, 221, 437}
\idxsub{his Foundations of Geometry}{437}
\idxmain{Historical notes}{6, 183, 202, 243, 321, 384}
\idxmain{Hypotenuse}{229}
\idxmain{Hypothesis}{42}
\idxmain{Implication}{41}
\idxmain{Incommensurable segments}{242}
\idxmain{Inequalities}{}
\idxsub{for angles}{155}
\idxsub{for numbers}{109}
\idxsub{for segments}{97, 116}
\idxmain{Infinite sequences}{}
\idxsub{of numbers}{302}
\idxsub{of polygonal paths inscribed in an arc}{323}
\idxmain{Infinite sum}{108}
\idxmain{Integers}{106}
\idxmain{Intersect}{64}
\idxmain{Irrational numbers}{107}
\idxmain{Isosceles triangle}{132}
\idxsub{base angles of}{132}
\idxmain{Jourdain, Philip}{437}
\idxmain{Leibniz}{7}
\idxmain{Le Verrier}{413}
\idxmain{Limit(s)}{}
\idxsub{area of rectangle as}{275}
\idxsub{circumference of circle as}{306}
\idxsub{length of circular arc as}{321}
\idxsub{of a sequence}{303}
\idxmain{Line(s)}{}
\idxsub{abstract character of}{14}
\idxsub{equations of}{394}
\idxsub{existence of parallels to}{199}
\idxsub{parallel}{198}
\idxsub{perpendicular}{158}
\idxsub{perpendicular to a plane}{366}
\idxsub{proof of straightness of}{177}
\idxsub{skew}{368}
\idxsub{straightness of}{63}
\idxsub{transversal of}{207}
\idxmain{Lobachevski}{7}
\idxmain{Locus}{346}
\idxsub{characterization of theorems on}{347}
\idxsub{cycloid as}{351}
\idxsub{definition of}{351}
\idxsub{hypocycloid as}{351}
\idxmain{Logic}{}
\idxsub{Aristotelian}{51}
\idxsub{deductive reasoning}{42}
\idxsub{definition of}{37}
\idxsub{drawing conclusions}{42}
\idxmain{Mathematics}{}
\idxsub{applied}{414}
\idxsub{consistency of}{438}
\idxsub{discovery or invention of}{438}
\idxsub{independent of experiments}{411}
\idxsub{models}{413}
\idxsub{pure}{414}
\idxsub{related to life}{438}
\idxsub{related to truth}{438}
\idxsub{research in}{438}
\idxsub{usefulness of}{413}
\idxmain{Measurement}{}
\idxsub{of angles in degrees}{330}
\idxsub{of angles in radians}{336}
\idxsub{of an angle inscribed in a circle}{338}
\idxsub{of arc length}{322}
\idxsub{of arcs in degrees}{334}
\idxsub{change of scale in}{118}
\idxsub{of segments}{110}
\idxsub{unit of length for}{110}
\idxmain{Midpoint}{}
\idxsub{of arc}{319}
\idxsub{formula for coordinates of}{403}
\idxsub{of segment}{133}
\idxmain{Newton, Isaac}{5}
\idxsub{his Principia}{5}
\idxmain{n-gon}{219}
\idxmain{Nonagon}{216}
\idxmain{Non-Euclidean geometry}{205}
\idxmain{Numbers}{}
\idxsub{irrational}{107}
\idxsub{positive integers}{106}
\idxsub{positive rational}{106}
\idxsub{ratio of integers}{106}
\idxsub{real}{108}
\idxsub{repeating decimals}{107}
\idxsub{role in Greek geometry}{102}
\idxsub{terminating decimals}{107}
\idxmain{Octagon}{222}
\idxmain{Ordinate}{388}
\idxmain{Orthocenter}{355}
\idxmain{Parallelogram}{223}
\idxsub{altitude of}{278}
\idxsub{base of}{279}
\idxmain{Pentagon}{218}
\idxmain{Perpendicular}{160}
\idxmain{Philosophy,}{}
\idxsub{epistemology}{408}
\idxsub{metaphysics}{408}
\idxsub{moral}{408}
\idxsub{political}{408}
\idxsub{religious}{408}
\idxsub{theology}{408}
\idxmain{Physics}{}
\idxsub{assumptions of}{412}
\idxsub{related to mathematics}{412}
\idxmain{Pi}{}
\idxsub{approximations for}{308}
\idxsub{calculation of rational approximations}{310}
\idxsub{definition of}{308}
\idxmain{Plane(s)}{}
\idxsub{abstract character of}{14}
\idxsub{left half of}{391}
\idxsub{lower half of}{391}
\idxsub{parallel}{369}
\idxsub{perpendicular to a line}{367}
\idxsub{right half of}{391}
\idxsub{upper half of}{391}
\idxmain{Plato}{6}
\idxmain{Point(s)}{}
\idxsub{abstract character of}{14}
\idxsub{collinear}{72}
\idxsub{of tangency}{294}
\idxsub{sides of on a line}{76}
\idxmain{Polygon(s)}{}
\idxsub{adjacent sides of}{215}
\idxsub{adjacent vertices of}{215}
\idxsub{apothem of regular}{298}
\idxsub{center of regular}{298}
\idxsub{central angle of regular}{298}
\idxsub{circumscribing a circle}{297}
\idxsub{classification of}{219}
\idxsub{consecutive angles of}{219}
\idxsub{convex}{218}
\idxsub{diagonal of}{219}
\idxsub{equiangular}{217}
\idxsub{equilateral}{219}
\idxsub{inscribed in a circle}{291}
\idxsub{interior and exterior of}{218}
\idxsub{radius of regular}{298}
\idxsub{regular}{218}
\idxsub{similar convex}{244}
\idxsub{simple}{216}
\idxsub{sum of exterior angles}{218}
\idxsub{sum of interior angles}{222}
\idxmain{Polygonal cell}{373}
\idxmain{Polygonal path}{}
\idxsub{definition of}{178}
\idxsub{endpoints of}{178}
\idxsub{inscribed in an arc}{322}
\idxsub{length of}{179}
\idxsub{sides of}{215}
\idxsub{vertices of}{178}
\idxmain{Positive axis}{386}
\idxmain{Positive direction}{386}
\idxmain{Postulate(s) (axioms, assumptions)}{62}
\idxsub{of Archimedes}{111}
\idxsub{of Group I (incidence)}{67}
\idxsub{of Group II (betweenness)}{69}
\idxsub{of Group III (linear congruence)}{93, 95}
\idxsub{of Group IV (angle congruence)}{121, 123}
\idxsub{necessity for}{102}
\idxsub{parallel}{202}
\idxmain{Proof}{39}
\idxsub{in algebra}{59}
\idxsub{independent of geometric figures}{63}
\idxsub{related to implication}{48}
\idxsub{relation of definitions to}{99}
\idxsub{techniques of}{51}
\idxmain{Properties, distinguished from relations}{127}
\idxmain{Proportion(al)}{237}
\idxsub{mean}{256}
\idxsub{sets of numbers}{239}
\idxsub{sets of segments}{239}
\idxmain{Ptolemy}{6}
\idxmain{Pythagoras}{3}
\idxmain{Pythagorean theorem}{257}
\idxsub{area proof of}{285}
\idxsub{converse of}{258}
\idxmain{Quadrants}{390}
\idxmain{Quadrilateral}{219}
\idxmain{Radian measure}{336}
\idxmain{Ratio}{237}
\idxsub{of integers}{106}
\idxsub{of segments}{239}
\idxmain{Rays}{77}
\idxsub{with same direction}{199}
\idxmain{Rectangle}{222}
\idxsub{area of}{273}
\idxsub{existence of}{224}
\idxmain{Relations}{127}
\idxmain{Rhombus}{222}
\end{multicols}
\end{document}
